\documentclass[12pt,a4paper]{article}

\usepackage[margin=1in]{geometry}
\usepackage[T1]{fontenc}
\usepackage[utf8]{inputenc}
\usepackage{lmodern}
\usepackage{enumitem}
\usepackage{xcolor}
\usepackage{hyperref}

\hypersetup{
	colorlinks=true,
	linkcolor=black,
	urlcolor=blue
}

\title{StoryTelling for the F1 DW Project}
\author{Giuseppe Rudi}
\date{ \today}

\begin{document}
	
	\maketitle
	
	\section*{Project idea}
	
	\begin{quote}
		\textit{Study Formula 1 competitive performance across the 2021 and 2022 seasons, in order to understand how the regulatory change affected team balance.}
	\end{quote}
	
	\section*{Storytelling}
	
	The final storytelling focuses on a more specific and interpretable subset of the domain.
	
	In particular, the final analysis considers only the three top teams:
	\begin{itemize}
		\item Ferrari
		\item Red Bull
		\item Mercedes
	\end{itemize}
	
	The storytelling is centered on comparing the behaviour and performance of their cars on a selected set of circuits, chosen according to their technical characteristics. For example, the selected circuits may represent different track profiles such as:
	\begin{itemize}
		\item high aerodynamic downforce circuits;
		\item power-sensitive circuits;
		\item tracks with slow corners;
		\item street circuits.
	\end{itemize}
	
	The analyses consider only dry-track conditions, excluding wet sessions, because in those situations the number of relevant factors becomes too large.
	
	
	The project does not consider the full telemetry at the finest level of detail. Instead, the analysis stops at a more balanced level of granularity, up to the three sectors that compose a lap.
	
	The analysis includes both:
	\begin{itemize}
		\item qualifying sessions;
		\item race sessions.
	\end{itemize}
	
	Free practice sessions are excluded, since they are less suitable for a clean and comparable performance analysis.
	
	Another important aspect of the storytelling is the comparison of car behaviour under different tyre compounds:
	\begin{itemize}
		\item Soft
		\item Medium
		\item Hard
	\end{itemize}
	
	Some situations are deliberately excluded from the final analysis. In particular, the project does not focus on:
	\begin{itemize}
		\item clear driver mistakes in a specific lap;
		\item laps or sessions performed in non-clean track conditions;
		\item penalties related to drivers, cars, or pit stops.
	\end{itemize}
	

	\section*{Details}
	
	\subsection*{Analytical idea}
	

	\begin{quote}
		\textit{Analyze how team performance changed between 2021 and 2022 after the regulatory shift.}
	\end{quote}

	\subsection*{Business Analyses. }
	
	
	\textbf{Main business question:}
	\begin{quote}
		\textit{How did the competitive balance among Ferrari, Red Bull, and Mercedes change from the 2021 season to the 2022 season after the regulatory change?}
	\end{quote}
	
	\vspace{2em}
    \textbf{Secondary questions:}
	\begin{itemize}
		\item How did the relative performance of Ferrari, Red Bull, and Mercedes change in qualifying and in race sessions?
		\item How did the performance of these three teams vary according to tyre compound (Soft, Medium, Hard)?
		\item On which types of circuits did the differences among these teams become more visible?
		\item Did the competitive gaps emerge more clearly in specific sectors of the lap?
		\item Did some teams perform better on high-downforce circuits, power-sensitive tracks, slow-corner circuits, or street circuits?
		\item How did car behaviour and performance change under different track and circuit characteristics?
		\item Did reliability-related aspects contribute to the competitive changes observed between 2021 and 2022?
	\end{itemize}
	
	
	
	\subsection*{Source Data}
	
	The source level of this project is not based on FastF1 only.  
	A secondary integrated source is also used for master data, namely the Ergast/Jolpica source accessed through the FastF1 \texttt{Ergast} interface.
	
	This choice is useful because FastF1 is mainly session-centered.  

	For this reason, the secondary Ergast/Jolpica source is used to build the following reconciled tables:
	\begin{itemize}
		\item \texttt{Driver}
		\item \texttt{Team}
	\end{itemize}
	

	

	The following data are extracted from FastF1:
	\begin{itemize}
		\item event schedule and race weekend information;
		\item session information for Qualifying and Race;
		\item session results;
		\item lap-level data;
		\item sector times for each lap;
		\item weather data;
		\item track status data.
	\end{itemize}
	
	In this way, the source level is divided into two complementary parts:
	\begin{itemize}
		\item Ergast/Jolpica for stable reference entities;
		\item FastF1 for detailed session and lap data.
	\end{itemize}
	
	
	\subsection*{Classification of Grand Prix}
	
	In order to answer some of the secondary business questions, the project needs a technical classification of both Grand Prix circuits and lap sectors.
	
	However, this type of classification is not directly available in the main adopted sources.  

	For this reason, the project introduces a set of \textbf{derived analytical attributes}.  
	These attributes are not copied directly from the source data.  
	Instead, they are obtained through a rule-based classification process supported by domain knowledge.

	\subsubsection*{Modeling choice}
	
	In principle, a more general and reusable design could separate the notion of circuit from the notion of circuit layout or circuit profile.  
	This would be useful in broader projects, because the same circuit may adopt different layouts in different years or racing categories.
	
	However, in this project the analysis is explicitly limited to the Formula 1 seasons 2021 and 2022.  
	Within this scope, the relevant circuit layouts are known and stable, and the analysis does not require modeling possible future layout changes.
	
	
	The circuit and its analytical classification are stored in a single table.  
	This table contains:
	\begin{itemize}
		\item the general descriptive information about the circuit;
		\item one categorical attribute describing the circuit as a whole;
		\item three categorical attributes describing the dominant technical nature of Sector 1, Sector 2, and Sector 3.
	\end{itemize}
	

	
	\subsubsection*{Global circuit category}
	
	Each circuit is assigned one global category among the following four classes:
	\begin{itemize}
		\item \texttt{Street}
		\item \texttt{PowerSensitive}
		\item \texttt{AeroSensitive}
		\item \texttt{Mixed}
	\end{itemize}
	
	Their meaning is the following:
	\begin{itemize}
		\item \texttt{Street}: urban or semi-urban circuit with barriers close to the track, lower average speed, and high importance of traction and precision;
		
		\item \texttt{PowerSensitive}: circuit where long straights, acceleration, top speed, and engine efficiency are dominant factors;
		
		\item \texttt{AeroSensitive}: circuit where aerodynamic load, stability in medium/high-speed corners, and aerodynamic efficiency are dominant factors;
		
		\item \texttt{Mixed}: circuit with a balanced profile, where no single technical characteristic clearly dominates.
		
	\end{itemize}
	
	\subsubsection*{Sector category}
	
	In addition to the global circuit category, each of the three sectors is assigned one category among the following four classes:
	\begin{itemize}
		\item \texttt{Power}
		\item \texttt{FastCorners}
		\item \texttt{SlowCorners}
		\item \texttt{Technical}
	\end{itemize}
	
	Their meaning is the following:
	\begin{itemize}
		\item \texttt{Power}: sector dominated by acceleration zones, long full-throttle sections, and heavy braking after high speed;
		
		\item \texttt{FastCorners}: sector dominated by medium/high-speed corners, where aerodynamic stability and cornering speed are especially important;
		
		\item \texttt{SlowCorners}: sector dominated by low-speed corners, traction, rotation, and corner exit performance;
		
		\item \texttt{Technical}: sector characterized by complex sequences of corners and direction changes, where precision and car balance are especially relevant.
		
	\end{itemize}
	
	

	
	\subsubsection*{Sources used }
	
	The list of Formula 1 race weekends considered in this project was derived from the official Formula 1 calendars for seasons 2021 and 2022:
	
	\begin{itemize}
		\item \url{https://www.formula1.com/en/racing/2021}
		\item \url{https://www.formula1.com/en/racing/2022}
	\end{itemize}
	
	These official pages were used to identify the race weekends effectively included in the two seasons and, therefore, the set of circuits relevant for the analysis. 
	
	The technical classification of the circuits was then based on descriptive circuit information collected from:
	
	\begin{itemize}
		\item \url{https://formula-timer.com/circuit}
	\end{itemize}
	
	This source was used mainly to observe the general technical profile of each circuit.
	
	For some sector-level interpretations, additional support was obtained from track-guide pages that describe the structure of specific sectors in detail. An example is:
	
	\begin{itemize}
		\item \url{https://www.tracktitan.io/post/monza-track-guide-sector-1-f1-2021}
	\end{itemize}
	
	These pages were not used as official structured datasets, but as descriptive support to assign the dominant category of each sector in a transparent and documented way. 


	
\end{document}